\centerline{\bf \Huge Метод от противного. Принцип Дирихле.}
\title{Метод от противного. Принцип Дирихле.}

\problem{1}Пятеро молодых рабочих получили на всех зарплату - 1500 рублей. Каждый из них хочет купить себе магнитофон ценой 320 рублей. Докажите, что кому-то из них придется подождать с покупкой до следующей зарплаты.

\problem{2.1}Докажите, что в любой футбольной команде есть два игрока, которые родились в один и тот же день недели.

\problem{2.2} Докажите, что среди жителей Москвы найдутся десять тысяч, празднующих день рождения в один и тот же день.

\problem{3} В лесу растет миллион елок. Известно, что на каждой из них не более 600000 иголок. Докажите, что в лесу найдутся две елки с одинаковым числом иголок.

\problem{4} В поход пошли 20 туристов. Самому старшему из них 35 лет, а самому младшему 20 лет. Верно ли, что среди туристов есть одногодки?

\problem{5} Докажите, что найдутся двадцать москвичей, имеющих одинаковое число волос на голове.(Известно, что у человека на голове не более 400000 волос, а в Москве не менее 8 миллионов жителей.)

\problem{6} В ковре размером 4 х 4 метра моль проела 15 дырок. Всегда ли можно вырезать коврик размером 1х1, не содержащий внутри дырок?

\problem{7} Можно ли разложить 44 шарика на 9 кучек так, чтобы количество шариков в разных кучках было различным?

\problem{8} В классе учатся 38 человек. Докажите, что среди них найдутся четверо, родившихся в один месяц.

\problem{9} Доказать, что если 21 человек собрали 200 орехов, то есть два человека, собравшие поровну орехов.

\problem{10} Какое наибольшее число королей можно поставить на шахматной доске так, чтобы никакие два из них не били друг друга?

\problem{11} Какое самое большое число ладей можно поставить на шахматную доску 8 на 8 так, чтобы они не били друг друга?

\problem{12} В клетках таблицы 3×3 расставлены числа –1, 0, 1.
Докажите, что какие-то две из восьми сумм по всем строкам, всем столбцам и двум главным диагоналям будут равны.